\subsection{EXPONENTIAL MAPPINGS}

THEOREM 4. Let G be a Lie group germ, L its Lie algebra, V an open neighbourhood of 0 in L, \(\phi\) an analytic mapping of V into G such that \(\phi(0) = 0\) and \(\mathrm{T}_0(\phi) = \mathrm{Id}_L\) . The following conditions are equivalent:

\begin{enumerate}
\def\labelenumi{(\roman{enumi})}
\item
  For all \(b \in \mathbf{L}\) , \(\phi((\lambda + \lambda')b) = \phi(\lambda b)\phi(\lambda'b)\) for \(|\lambda|\) and \(|\lambda'|\) sufficiently small.
\item
  For all \(b \in \mathbf{L}\) and every integer \(n > 0\) , \(\phi_{*}(b^{n})\) is homogeneous of degree \(n\) in

\(\mathrm{U}(\mathrm{G})\) (\(\mathrm{T}_{0}^{(\infty)}(\mathrm{L})\) is identified with \(\mathrm{TS}(\mathrm{L})\) and \(b^{n}\) is evaluated in \(\mathrm{TS}(\mathrm{L})\)).

\item
  The mapping \(\phi_{*}\) of TS(L) into U(G) is compatible with the graduations of TS(L) and U(G).
\item
  The mapping \(\phi_*\) of TS(L) into U(G) is the canonical mapping of TS(L) into the enveloping algebra of L.
\item
  There exist a norm on \(\mathbf{L}\) defining the topology of \(\mathbf{L}\) and such that
\end{enumerate}

\[
\| [ x, y ] \| \leqslant \| x \| \| y \|
\]

for all \(x, y\) in \(\mathbf{L}\) and an open subgroup germ \(\mathbf{W} \subset \mathbf{V}\) of the Lie group germ defined by \(\mathbf{L}\) (no. 2) such that \(\phi | \mathbf{W}\) is an isomorphism of \(\mathbf{W}\) onto an open Lie subgroup germ of \(\mathbf{G}\) .

(v) \(\Rightarrow\) (i): obvious, for \((\lambda b).(\lambda^{\prime}b) = (\lambda +\lambda^{\prime})b\) in W for \(|\lambda |\) and \(|\lambda '|\) sufficiently small.
(i) \(\Rightarrow\) (ii): suppose that condition (i) is satisfied. Let \(b \in L\) . Let \(\psi\) be the restriction of \(\phi\) to \(V \cap Kb\) . By hypothesis, there exists a symmetric neighbourhood \(T\) of 0 in the additive Lie group \(Kb\) such that \(\psi|T\) is a morphism of the Lie group germ \(T\) into \(G\) . Hence

\[
\phi_ {*} (b ^ {n}) = (\psi | \mathrm{T}) _ {*} (b ^ {n}) = ((\psi | \mathrm{T}) _ {*} (b)) ^ {n} = (\phi_ {*} (b)) ^ {n},
\]

so that \(\phi_{*}(b^{n})\) is homogeneous of degree \(n\) in U(G).

(ii) \(\Rightarrow\) (iii): this follows from the fact that \(\mathsf{TS}^n (\mathbf{L})\) is the vector subspace of \(\mathsf{TS}(\mathbf{L})\) generated by the \(n\) -th powers of the elements of \(\mathbf{L}\) (Algebra, Chapter IV, § 5, Proposition 5).
(iii) \(\Rightarrow\) (iv): the canonical mapping of TS(L) into the enveloping algebra of L is the unique morphism of graded cogebras mapping 1 to 1 and extending \(\mathrm{Id}_{\mathrm{L}}\) (Chapter II, § 1, no. 5, Remark 3). Now \(\phi_{*}\) is a cogebra morphism and \(\phi_{*}|_{\mathrm{L}} = \mathrm{Id}_{\mathrm{L}}\) by hypothesis. If condition (iii) holds, it is seen that condition (iv) also holds.
(iv) \(\Rightarrow\) (v): suppose that condition (iv) is satisfied. Choose a norm on L defining the topology of L and such that \(\| [x,y]\| \leqslant \| x\| \| y\|\) for all \(x,y\) in L. Let H be the Lie group germ defined by the normed Lie algebra L. By Theorem 1, there exist an open subgroup germ \(S\subset V\) of H and an isomorphism \(\phi^{\prime}\) of S onto an open subgroup germ of G. As we know already that (v) \(\Rightarrow\) (iv), the mapping \(\phi_{*}\) of TS(L) into U(G) is the canonical mapping of TS(L) into the enveloping algebra of L. Thus \(\phi_{*}(t) = \phi_{*}^{\prime}(t)\) for all \(t\in T_0^{(\infty)}(L)\) . As \(\phi\) and \(\phi^{\prime}\) are analytic, \(\phi\) and \(\phi^{\prime}\) coincide on a neighbourhood of 0.

DEFINITION 1. Let G be a Lie group germ and L its Lie algebra. An exponential mapping of G is any analytic mapping \(\phi\) defined on an open neighbourhood of 0 in L, with values in G and satisfying the conditions of Theorem 4.

Theorem 4 implies immediately that, for every Lie group germ G, there exists an exponential mapping of G and that two exponential mappings of G coincide on a neighbourhood of 0.

Examples. (1) Let G be the additive group of a complete normable space E. The canonical isomorphism of L(G) onto E satisfies condition (i) of Theorem 4 and is therefore an exponential mapping of G.

\begin{enumerate}
\def\labelenumi{(\arabic{enumi})}
\setcounter{enumi}{1}
\tightlist
\item
  Let A be a complete normed unital associative algebra. Let \(\mathbf{A}^*\) be the Lie group consisting of the invertible elements of A. We identify \(\mathrm{L}(\mathrm{A}^*)\) with A (§3, no. 9, Corollary to Proposition 33). If \(\mathbf{K} = \mathbf{R}\) or \(\mathbf{C}\) , we know that the mapping exp of A into \(\mathbf{A}^*\) defined in Chapter II, §7, no. 3 satisfies condition (i) of Theorem 4 and hence is an exponential mapping. Now let K be ultrametric. Let \(p\) be the characteristic of the residue field of K. If \(p \neq 0\) , let \(\lambda = |p|^{1/(p-1)}\) ; if \(p = 0\) , let \(\lambda = 1\) . Let U be the set of \(x \in \mathbf{A}\) such that \(\|x\| < \lambda\) . We know (Chapter II, §8, no. 4) that the mapping exp of U into \(\mathbf{A}^*\) satisfies condition (i) of Theorem 4 and hence is an exponential mapping. Note that U is an additive subgroup of A.
\end{enumerate}

This example explains the terminology adopted in Definition 1.

Let \(G\) be a Lie group germ and \(\phi\) an exponential mapping of \(G\) . Then \(\phi\) is étale at 0 and hence there exists an open neighbourhood \(U\) of 0 in \(L(G)\) such that \(\phi(U)\) is open in \(G\) and \(\phi|U\) is an isomorphism of the analytic manifold \(U\) onto the analytic manifold \(\phi(U)\) .

A canonical chart (of the first species) on G is a chart \(\psi\) on the analytic manifold G whose inverse mapping is an exponential mapping. If further G is finite-dimensional and a basis of L(G) is chosen, the coordinate system defined by \(\psi\) and this basis in the domain of \(\psi\) is called a canonical coordinate system (of the first species).

PROPOSITION 3. Let \(\mathbf{G}\) be a Lie group germ, \(\mathbf{L}\) its Lie algebra and \(\phi\) an exponential mapping of \(\mathbf{G}\) . Let \(\mathbf{L}_1, \ldots, \mathbf{L}_n\) be vector subspaces of \(\mathbf{L}\) such that \(\mathbf{L}\) is the topological direct sum of \(\mathbf{L}_1, \ldots, \mathbf{L}_n\) . The mapping

\[
(b _ {1}, b _ {2}, \dots , b _ {n}) \mapsto \theta (b _ {1}, b _ {2}, \dots , b _ {n}) = \phi (b _ {1}) \phi (b _ {2}) \dots \phi (b _ {n}),
\]

defined on an open subset of \(L_{1} \times L_{2} \times \cdots \times L_{n}\) , is analytic. The tangent mapping at \((0, 0, \ldots, 0)\) to \(\theta\) is the canonical mapping of \(L_{1} \times \cdots \times L_{n}\) into L.

Let \(k_{i}\) be the canonical injection of \(L_{i}\) into \(L_{1} \times L_{2} \times \cdots \times L_{n}\) . Then, for all \(b \in L_{i}\) , \((\mathrm{T}_{(0,\ldots,0)}\theta)(\mathrm{T}_{0}\phi_{i})(b) = (\mathrm{T}_{0}\phi)(b) = b\) and hence \((\mathrm{T}_{(0,\ldots,0)}\theta)|L_{i}\) is the canonical injection of \(L_{i}\) into L.

In particular, \(\theta\) is étale at \((0, 0, \ldots, 0)\) . Its restriction to a sufficiently small open neighbourhood U of \((0,0,\ldots,0)\) has an open image in G and is an isomorphism of the manifold U onto the manifold \(\theta(\mathrm{U})\) . The inverse mapping \(\eta\) of \(\theta(\mathrm{U})\) onto U is called a canonical chart of the second species of G, associated with the given decomposition of L as a direct sum. If further G is finite-dimensional and each \(L_{i}\) is generated by a non-zero vector \(e_{i}\) , the coordinate system on \(\theta(\mathrm{U})\) defined by \(\eta\) and the \(e_{i}\) is called a canonical coordinate system of the second species.

PROPOSITION 4. Let G be a Lie group germ and \(\phi\) an injective exponential mapping of G. For all \(x, y\) in L(G),


\[
x + y = \lim _ {\lambda \in \mathbb {K} ^ {*}, \lambda \rightarrow 0} \lambda^ {- 1} \phi^ {- 1} (\phi (\lambda x) \phi (\lambda y))\tag{1}
\]

\[
[ x, y ] = \lim _ {\lambda \in \mathbb {K} ^ {*}, \lambda \rightarrow 0} \lambda^ {- 2} \phi^ {- 1} (\phi (\lambda x) \phi (\lambda y) \phi (- \lambda x) \phi (- \lambda y))\tag{2}
\]

(note that \(\phi^{-1}(\phi(\lambda x)\phi(\lambda y))\) and \(\phi^{-1}(\phi(\lambda x)\phi(\lambda y)\phi(-\lambda x)\phi(-\lambda y))\) are defined for \(|\lambda|\) sufficiently small).

Let \(\mathbf{L} = \mathbf{L}(\mathbf{G})\) be given a norm defining the topology of \(\mathbf{L}\) and such that \(\| [x,y] \| \leqslant \| x \| \| y \|\) for all \(x, y\) in \(\mathbf{L}\) . Using Theorems 2 and 4, it can be assumed that \(\mathbf{G}\) is the Lie group germ defined by \(\mathbf{L}\) and that \(\phi = \mathrm{Id}_{\mathbf{G}}\) . Let \((x,y) \mapsto x.y\) denote the product in the group \(\mathbf{G}\) . The formulae to be proved can then be written


\[
x + y = \lim _ {\lambda \in \mathbf {K} ^ {*}, \lambda \rightarrow 0} \lambda^ {- 1} ((\lambda x). (\lambda y))\tag{3}
\]

\[
[ x, y ] = \lim _ {\lambda \in \mathbb {K} ^ {*}, \lambda \rightarrow 0} \lambda^ {- 2} ((\lambda x). (\lambda y). (- \lambda x). (- \lambda y)).\tag{4}
\]

There exists an open neighbourhood V of 0 in K such that the function

\[
\lambda \mapsto f (\lambda) = (\lambda x). (\lambda y)
\]

is defined and analytic on V. By Chapter II, § 6, no. 4, Remark 2, the expansion of \(f\) as an integral series about the origin is

\[
\lambda (x + y) + \frac {1}{2} \lambda^ {2} [ x, y ] + \dots
\]

and this proves (3). On the other hand, for \(u, v\) in G and \(\| u \|\) , \(\| v \|\) sufficiently small, \(u.v\) is an analytic function of \((u, v)\) and the terms of degrees 1 and 2 in the expansion of this function as an integral series about the origin are \(u + v + \frac{1}{4}[u, v]\) . By Differentiable and Analytic Manifolds, R, 3.2.7 and 4.2.3, the terms of degrees 1 and 2 in the expansion of the function \(f(\lambda)f(-\lambda)\) as an integral series about the origin are the terms of degrees 1 and 2 in

\[
f (\lambda) + f (- \lambda) + \frac {1}{2} [ f (\lambda), f (- \lambda) ]
\]

or also in

\[
\begin{array}{r l} \lambda (x + y) + \frac {1}{2} \lambda^ {2} [ x, y ] - \lambda (x + y) + \frac {1}{2} \lambda^ {2} [ x, y ] & \\ + \frac {1}{2} [ \lambda (x + y), - \lambda (x + y) ] = \lambda^ {2} [ x, y ] \end{array}
\]

and this proves (4).

PROPOSITION 5. Let G be a Lie group, k a non-discrete closed subfield of K, G' the group G considered as a Lie group over k and \(\phi\) (resp. \(\phi'\) ) an exponential mapping of G (resp. G'). Then \(\phi\) and \(\phi'\) coincide on a neighbourhood of 0.

\(\phi\) satisfies hypothesis (i) of Theorem 4 relative to \(G'\) and is therefore an exponential mapping of \(G'\) .

PROPOSITION 6. Let G be a Lie group germ, L its Lie algebra and \(\phi: V \to G\) an exponential mapping of G. For all \(x \in V\) , let \(\mathrm{T}_x(L)\) be identified with L, so that the right differential \(\varpi(x)\) of \(\phi\) at \(x\) is a linear mapping of L into L. For \(x\) sufficiently close to 0,

\[
\varpi (x) = \sum_ {n \geqslant 0} \frac {1}{(n + 1) !} (\mathrm{ad} x) ^ {n}.
\]

Let \(\mathbf{L}\) be given a norm compatible with its topology and such that \(\| [x,y]\| \leqslant \| x\| \| y\|\) for all \(x,y\in L\) . It suffices to consider the case where \(G\) is the Lie group germ defined by \(L\) and \(\phi = \mathrm{Id}_{\mathrm{G}}\) . By definition, \(\varpi (x)\) is then the tangent mapping at \(x\) to the mapping \(y\mapsto y.x^{-1}\) of \(G\) into \(G\) . If \(H(X,Y)\) denotes the Hausdorff series, \(\varpi (x)\) is therefore, for \(\| x\|\) sufficiently small, the tangent mapping at 0 to the mapping \(y\mapsto H(x + y, - x)\) of \(G\) into \(G\) . In \(H(X + Y, - X)\) , the sum of terms of the first degree in \(Y\) is

\[
\sum_ {m \geqslant 0} \frac {1}{(m + 1) !} (\mathrm{adX}) ^ {m} \mathrm{Y}
\]

(Chapter II, § 6, no. 5, Proposition 5). The proposition then follows from Differentiable and Analytic Manifolds, R, 3.2.4 and 4.2.3.

Let G be a Lie group germ and \(t \in K\) . A t-th power mapping of G is any mapping, defined and analytic on an open neighbourhood of e, with values in G and coinciding on a neighbourhood of e with a mapping

\[
g \mapsto \phi (t \phi^ {- 1} (g))
\]

where \(\phi\) is an injective exponential mapping of G.

PROPOSITION 7. (i) If \(t \in \mathbf{Z}\) , a \(t\) -th power mapping coincides on a neighbourhood of \(e\) with the mapping \(g \mapsto g^t\) .

\begin{enumerate}
\def\labelenumi{(\roman{enumi})}
\setcounter{enumi}{1}
\item
  The tangent mapping at \(e\) to a \(t\) -th power mapping is the homothety of ratio \(t\) .
\item
  If \(h\) is a \(t\) -th power mapping and \(h'\) a \(t'\) -th power mapping of \(G\) , \(h \circ h'\) is a \((tt')\) -th power mapping and \(g \mapsto h(g)h'(g)\) is a \((t + t')\) -th power mapping.
\item
  If \(h\) is a \(t\) -th power mapping and \(u \in \mathbf{U}^n(\mathbf{G})\) , then \(h_*(u) = t^n u\) .
\end{enumerate}

It suffices to prove the proposition when G is the Lie group germ defined by a complete normed Lie algebra and when the t-th power mappings considered are constructed using the exponential mapping \(\phi = \mathrm{Id}_{\mathrm{G}}\) . But in that case everything is obvious.

